\scanpage{182}
% L01–05
deux inclusions (8) (9), inverses l'une de l'autre, implique l'autre i.e. l'égalité (10), i.e. de nombreux cas où l'inclusion 3) est en fait une égalité
\[
Zs_0=S^{G_{s_0}}.\tag*{(15)}
\]

% L06–13
Voici des cas où il en est ainsi
\[
\left\{\begin{aligned}
\text{a)}\ &S\text{ fini}\quad(\text{car tout hom surjectif de }S\text{ sur}\\[-2pt]
&\hspace{7em}\text{lui }=\text{ est alors bijectif})\\
\text{b)}\ &\deleted{\text{les}}\ G_s\text{ finis}\quad(\text{raison directe pour les inclusions}\\
&\hspace{7em}\text{de la forme (8) }\unsure{(9)|(8)}\,)\\
\text{c)}\ &\deleted{\text{les}}\ G\text{ est un groupe de Lie réel, }\deleted{\text{et}}\\
&\text{les }G_s\text{ sont fermés et n'ont qu'un nb}\\
&\text{fini de composantes connexes.}
\end{aligned}\right.\tag*{(16)}
\]

% L14–16
\deleted{\unsure{On}} On peut appeler aussi points \inserted{$G$-}\uline{antipodiques} d'un $s_0\in S$ les points \deleted{$s\ne s_0$ \unsure{qui sont|tels}} \deleted{\unsure{$s$}} $\in S^{G_{s_0}}$, i.e. fixés par tout élément de $G$ qui \deleted{\unsure{le}} fixe $s_0$.

% L17–23
Dans les cas envisagés a) b) c), ce sont aussi les pts $G$-équivalents à $s_0$, qu'on peut aussi appeler \uline{strictement antipodiques} à $s_0$. On peut appeler application \inserted{$G$-}antipodique de $S$ \deleted{\unsure{toute}} application qui commute à $G$ (donc \unsure{néc\textsuperscript{t}} surjective) -- \deleted{\unsure{ainsi}} et application \inserted{$G$-}antipodique une \inserted{telle} application bijective.

% archive sheet number: 181
